\documentclass[11pt,a4paper]{article}
\usepackage[T1]{fontenc}
\usepackage[utf8]{inputenc}
\usepackage{amsmath,amssymb,amsthm,hyperref}
\usepackage[margin=27mm]{geometry}
\newtheorem{proposition}{Proposition}
\title{Une conséquence de la borne quantitative de crible pour le problème d'Erdős 687}
\author{Note de travail de la session de recherche de francoisbrunel avec Codex}
\date{28 novembre 2026}
\begin{document}
\maketitle

Pour chaque premier $p\le x$, on choisit une classe $a_p\bmod p$.
On note $Y(x)$ la plus grande longueur d'un intervalle $[1,y]$ que l'union
des classes peut recouvrir. L'objet de cette note est une borne supérieure.
L'ordre de grandeur demandé dans la partie (i) reste à déterminer.

\begin{proposition}
Le théorème 1.2 de \cite{survivors} entraîne l'existence de constantes
absolues $\eta>0$ et $C>0$ telles que, pour tout $x$ suffisamment grand,
\[
                  Y(x) \le C\,\frac{x^2}{(\log x)^{2+\eta}} .
\]
Les constantes de la preuve du théorème permettent de prendre $\eta=1/200$.
\end{proposition}

\begin{proof}
\emph{Entrée quantitative.}
Écrivons $N(z)=\lfloor z^2/(\log z)^2\rfloor$. Le théorème 1.2 de
\cite{survivors} et la formule de Mertens donnent, pour une constante
absolue $a>0$, pour tout $z$ suffisamment grand et pour tous les choix de
classes aux premiers $r\le z$,
\[
 |S| \ge a N(z) \frac{\log\log z}{(\log z)^2},
 \qquad S=\{1\le n\le N(z): n\not\equiv a_r\bmod r
                                      \ \text{pour tout premier }r\le z\}.
\]
En effet, la borne du théorème est $cN(z)V_0/B^2$, avec
$w=(\log z)/(\log\log z)^2$, $B=(\log z)/\log w$ et
$V_0=\prod_{r\le w}(1-1/r)\sim e^{-\gamma}/\log w$.
Ainsi $V_0/B^2\asymp(\log\log z)/(\log z)^2$.

\emph{Choix de l'échelle.}
Posons $L=\log x$, $z=x/L^\delta$ et $N=N(z)$, où $0<\delta\le 1$
sera fixé à la fin. Uniformément pour ces $\delta$, on a
$L-\log L\le\log z\le L$. Il existe donc $a_0>0$ absolue telle que
\begin{equation}\label{eq:stock}
                 |S| \ge a_0 N\frac{\log L}{L^2}.
\end{equation}

\emph{Contribution d'un premier supplémentaire.}
Pour $z<p\le x$, les entiers de $[1,N]$ congrus à $a_p\bmod p$
forment une progression de raison $p$, dont les indices forment un
intervalle de $M$ entiers, avec $|M-N/p|\le1$. Pour $x$ assez grand,
$x^{1/2}\le M\le x$ uniformément en $p$ et $\delta$.
En effet $N/x\gg x/L^4$ et $N/z\le z/(\log z)^2$.
Pour chaque premier $r\le M^{1/4}$, on a $r\le z<p$ et $p$ est
inversible modulo $r$. La condition $n\not\equiv a_r\bmod r$ exclut
donc une classe des indices modulo $r$.
Le crible linéaire supérieur appliqué à ces $M$ indices au paramètre
$u=4$ donne au plus $C_1 M/\log M$ survivants, uniformément dans
les classes. Pour retrouver la formulation par divisibilité, on utilise
le théorème chinois sur les classes d'indices puis on translate
l'intervalle. Voir \cite[formule (1)]{granville}. En particulier,
avec $C_2$ absolue,
\begin{equation}\label{eq:class}
        |\{n\in S:n\equiv a_p\bmod p\}|
                      \le C_2\,\frac{N}{pL}.
\end{equation}

\emph{Somme sur les premiers.}
Le majorant de Tchebychev $\pi(t)\ll t/\log t$ implique
$\sum_{u<p\le 2u}p^{-1}\ll1/\log u$.
On découpe $(z,x]$ en au plus $1+\log(x/z)/\log 2$ intervalles
de ce type, le dernier pouvant être tronqué. Comme
$\log(x/z)=\delta\log L$ et $\log z\ge L/2$, on obtient, avec une
constante absolue indépendante de $\delta$,
\begin{equation}\label{eq:harmonic}
            \sum_{z<p\le x}\frac1p \ll \frac{1+\delta\log L}{L}.
\end{equation}
Par la borne de l'union, \eqref{eq:class} et \eqref{eq:harmonic},
les classes supplémentaires retirent de $S$ au plus
$D N(1+\delta\log L)/L^2$ entiers, avec $D>0$ absolue.
Prenons $\delta=\min(1/2,a_0/(4D))$. Dès que
$D\le(a_0/4)\log L$, elles retirent au plus la moitié du minorant
\eqref{eq:stock}. Il reste un entier dans $[1,N]$, quels que soient
les restes initiaux. Par restriction de tout intervalle recouvert,
$Y(x)<N$. Or $N\asymp x^2/L^{2+2\delta}$. On conclut avec
$\eta=2\delta$.
\end{proof}

\paragraph{Vérification de l'échelle par une deuxième sommation.}
Pour $\delta$ fixé, la version du théorème des nombres premiers avec
reste classique, suivie d'une sommation partielle, donne
\[
 \sum_{x/L^\delta<p\le x} \frac1p
   = \log\!\left(\frac{L}{L-\delta\log L}\right)+o(1/L)
   \sim \delta\,\frac{\log L}{L}.
\]
Cette vérification retrouve le facteur $\delta$ indispensable dans la
comparaison des constantes.

\paragraph{Justification du choix explicite $\eta=1/200$.}
La proposition « Positive reference margin » de \cite{survivors}
donne $\lim B^2 P_{\rm e}/e^\gamma>0.02$ à la racine. On peut donc
utiliser $c_L=0.02 e^\gamma$ dans la comparaison de la section finale
de cette prépublication. Elle établit $|S|\ge(c_L/2)NV_0/B^2$.
Comme $\log w\sim\log L$ et $\log z\sim L$, ceci donne, uniformément
dans les classes pour $\delta$ fixé,
\[
   |S|\ge (0.01-o(1))N\log L/L^2 .
\]
La majoration de progression de Banks, Ford et Tao
\cite[section 2.3]{bft} s'applique déjà après le crible aux premiers
$r\le\sqrt{N/\log N}$, tous inférieurs à $z$. Pour chaque premier
$z<p\le x$ elle donne au plus $2(N/p)/\log(2N/p)$ entiers dans la
classe considérée, en utilisant l'intervalle $[0,N]$ qui contient
$[1,N]$. Les classes aux petits premiers sont arbitraires. Ici
$\log(2N/p)=L+O(\log L)$ uniformément en $p$, ce qui majore la somme
des pertes par $(2\delta+o(1))N\log L/L^2$. La même constante suit
du crible linéaire supérieur au paramètre $2$ sur les indices de
chaque progression, par $F(2)=e^\gamma$ et la formule de Mertens.
Avec $\delta=1/400$ le coefficient des pertes tend au plus vers
$0.005$, strictement inférieur à $0.01$. La conclusion suit avec
$2\delta=1/200$. La constante multiplicative et le seuil en $x$ ne
sont pas explicités.

\paragraph{Dérivation formelle de la constante du crible supérieur.}
Soit $R\le z$ un rayon entier, $G(R)$ la normalisation de Selberg pour
les premiers jusqu'à $z$, et $H_R=\sum_{n\le R}1/n$. Pour des classes
arbitraires et $J$ positions consécutives, la vérification donne
\[
     |S_z(J)|\le\frac{J}{\log(R+1)}+(R H_R)^2,\qquad R\ge1.
\]
En effet, groupons les entiers $n\le R$ par leur radical $d$.
La somme de leurs inverses est au plus
$\prod_{p\mid d}\sum_{e\ge1}p^{-e}=\prod_{p\mid d}1/(p-1)$,
le terme que $d$ apporte à $G(R)$. On obtient donc $G(R)\ge H_R$.
Pour borner le reste, les masses diagonales absolues $u_\ell$ du crible
optimisé vérifient $\sum u_\ell=1$ et sont nulles pour $\ell>R$.
L'inversion de Möbius donne
$|\lambda_d|\le d\sum_{d\mid\ell}u_\ell$. En sommant d'abord sur $d$,
on trouve $\sum_d|\lambda_d|\le\sum_\ell u_\ell\sigma(\ell)\le R H_R$,
car $\sigma(\ell)=\ell\sum_{d\mid\ell}1/d$. Le terme d'erreur du
crible est au plus le carré de cette somme. Pour $\varepsilon>0$, le
choix $R=\lfloor J^{1/(2+\varepsilon/2)}\rfloor$ donne, uniformément
pour $z\ge J$ et $J$ assez grand,
$|S_z(J)|\le (2+\varepsilon)J/\log J$. La preuve Lean utilise la
diagonalisation de Selberg du projet cité en \cite{survivors}.

\paragraph{Une limite de la majoration par progression.}
Pour un entier $M$ et $t$ premiers distincts $p_1,\ldots,p_t>z\ge M$,
posons $u=t\pi(M)<M$. Dans l'intervalle $[1,M\max_i p_i]$, il existe
des classes aux premiers $r\le z$ qui épargnent simultanément tous
les éléments de
\[
 T=\{p_iq:\ 1\le i\le t,\quad q\text{ premier},\quad u<q\le M\}.
\]
En effet, pour $r\le u$, prescrire la classe $0$ épargne tout $T$.
Pour $u<r\le z$, on a $|T|\le t\pi(M)=u<r$. Il existe donc une classe
modulo $r$ évitant $T$. L'unicité de la factorisation et $p_i>M$
rendent les $t$ sous-ensembles disjoints. En prenant ensuite la classe
$0$ pour chacun des $p_i$, ils retirent ensemble au moins
$t(\pi(M)-\pi(u))$ survivants. Si $t=o(\log M)$, alors $u=o(M)$
et $\pi(u)=o(\pi(M))$, par monotonie et le théorème des nombres
premiers. Lorsque les $p_i$ sont de tailles comparables, cette perte
est du même ordre $tM/\log M$ que la somme des majorants individuels.
Ce constat s'applique aussi à chacune des classes d'un groupe fixe,
avec le même choix de classes pour tous les petits premiers.

\paragraph{Pertes simultanées sur toute la bande.}
Écrivons $D(N;z,x)$ pour le maximum de $|S_z\setminus S_x|$ sur les choix
de classes, où $S_t$ est l'ensemble qui échappe aux premiers jusqu'à $t$
dans $[1,N]$. Pour $N=N(z)$, $z=x/L^\delta$ et $\delta>0$ fixé, posons
\[
 A=\frac{N}{L}\sum_{z<p\le x}\frac1p
   \sim \delta N\frac{\log L}{L^2}.
\]
Alors $(1-o(1))A\le D(N;z,x)\le(2+o(1))A$. La majoration est celle
déjà établie par le crible supérieur. Pour la minoration, prenons
$U=N/(xL)$, $V=N/z$ et protégeons les produits $n=pq\le N$ où
$z<p\le x$ et $U<q\le N/p$ sont premiers. Ces produits sont tous
distincts car $q\le V<z$. Le théorème des nombres premiers et
$U\le(N/p)/L$ donnent
\[
 |T|=\sum_{z<p\le x}\bigl(\pi(N/p)-\pi(U)\bigr)=(1+o(1))A,
\]
puisque $\log(N/p)\sim L$ uniformément. Jusqu'à $z$, fixons le reste
$0$ pour les premiers en dehors de $(U,V]$. Aucun ne divise un
élément de $T$. Pour chaque premier $r$ dans $(U,V]$, choisissons
indépendamment un reste uniforme \emph{non nul}. Un $n\in T$ survit
avec probabilité au moins
\[
       \prod_{U<r\le V}\left(1-\frac1{r-1}\right)=1-o(1).
\]
En effet, $\log U\sim L$ et $V/U=L^{1+\delta}$, et la somme des
$1/(r-1)$ sur cette bande est $O(\log L/L)$. Par moyenne, un choix
épargne au moins $(1-o(1))|T|$ de ces produits. Attribuer ensuite
le reste $0$ à tous les $p$ de $(z,x]$ les retire ensemble.

Pour $z=x^\beta$ avec $1/2<\beta<1$ fixé, le même argument donne
\[
              D(N;z,x)\asymp_\beta N/\log x .
\]
En effet, $\log U\sim(2\beta-1)L$, $\log V\sim\beta L$, et le
produit de survie tend vers $(2\beta-1)/\beta>0$ par la formule de
Mertens (remplacer $1/r$ par $1/(r-1)$ coûte $O(1/U)$).
L'estimation uniforme $\log(N/p)\asymp_\beta L$ et
$\sum_{z<p\le x}1/p\to\log(1/\beta)>0$ suffisent pour le comptage.
La majoration par progression fournit l'autre inégalité. Ainsi les
pertes maximales peuvent réellement dépasser d'un facteur divergent le
minorant garanti par le théorème 1.2 lorsque la bande a cette largeur.
Les configurations réalisant les deux extrêmes peuvent être différentes.
Une estimation reliant la perte à l'effectif initial pourrait encore
améliorer le résultat.

Une variante tient compte des intersections sur le même ensemble $S$
de survivants des petits premiers. Pour $p$ dans la bande restante $P$,
notons $C_p=\{n\in S:n\equiv a_p\bmod p\}$. Pour chaque $p$, choisissons
au plus un premier antérieur $\rho(p)\in P$, $\rho(p)<p$, et donnons la
classe vide aux parents absents. Alors
\[
 \left|\bigcup_{p\in P} C_p\right|
 \le \sum_{p\in P}\bigl(|C_p|-|C_p\cap C_{\rho(p)}|\bigr).
\]
En effet, pour tout entier couvert, sa première classe le compte encore
après le retrait de son intersection avec le parent. Cette majoration
sur une forêt est une forme de la borne de Hunter
(\href{https://www.cambridge.org/core/journals/journal-of-applied-probability/article/an-upper-bound-for-the-probability-of-a-union/092D711504BA968EF0D1D903A2685D60}{Journal of Applied Probability, 1976}).
Pour exclure les couvertures complètes à une longueur donnée, il suffit
donc que, pour chaque choix des résidus, il existe des parents pour
lesquels cette somme est strictement inférieure à $|S|$. Le choix des
parents peut dépendre des résidus. Pour conclure à une nouvelle borne
asymptotique, il reste à prouver ce critère uniformément en $x$ à la
longueur visée.

Cette correction a toutefois une limite uniforme particulièrement forte
dans la bande utilisée ici. Lorsque $H<pq$ pour chaque paire distincte
de premiers de $P$, deux classes distinctes par leur module ont au plus
un point commun dans $[1,H]$. Deux points communs distincts auraient
une différence positive multiple de $pq$ et strictement inférieure à
$pq$. Le gain total de la forêt est donc au plus $|P|$ pour tout choix
de résidus et de parents. À $z=x/(\log x)^{1/400}$ et
$H=\lfloor z^2/(\log z)^2\rfloor$, on a finalement $H<z^2$, tandis que
$|P|\le x=o(H\log\log x/(\log x)^2)$. La forêt seule ne peut donc
économiser une proportion fixe de l'unité de réserve à cette échelle.

Une autre condition permet d'utiliser toutes les paires, en tenant
compte des triples pour contrôler les recouvrements concentrés sur
quelques positions. Pour $i\in S$, soit $k_i$ le nombre de classes qui
le touchent. Posons $d_i=\max(k_i-1,0)$, $D=\sum_i d_i$,
$B=\sum_i\binom{k_i}{2}$ et $T=\sum_i\binom{k_i}{3}$.
L'inégalité de Cauchy--Schwarz pondérée par $d_i$ donne
\[
 (2B)^2=\Big(\sum_i k_id_i\Big)^2
 \le D\sum_i k_i^2d_i=D(6T+4B).
\]
Si $B>0$, le gain sur la somme brute des tailles des classes vaut
donc au moins $2B^2/(3T+2B)$. Lors d'une couverture complète,
$D=\sum_i k_i-|S|$ exactement. Cette identité donne un test
d'impossibilité sans division, y compris si $B=0$.
Pour en déduire un nouveau majorant de $Y$, il reste à établir des
estimations uniformes sur ces moments de multiplicité.

\paragraph{Une condition pour améliorer l'exposant de $x$.}
Pour préciser l'obstacle à un gain sur la puissance de $x$, fixons un
entier $m\ge2$ et écrivons $z=k^m$, $x=k^{m+1}$, $N=z^2$. Soit $S$
l'ensemble des positions de $[1,N]$ épargnées par les premiers jusqu'à $z$,
et $C_p$ la classe choisie modulo $p$, restreinte à $S$. Pour tout $k$
assez grand (à $m$ fixé) et tout choix des classes,
\[
 \sum_{z<p\le x}|C_p|\le\frac{24}{m}\frac{N}{\log z}.
\]
En effet $J=\lfloor N/p\rfloor+1$ vérifie
$k^{m-1}\le J\le z$ et $\log J\ge\frac12\log z$.
Le crible supérieur sur les indices de la progression donne
$|C_p|\le3J/\log J\le12N/(p\log z)$. La borne de Mertens utilisée
ci-dessus donne $\sum_{z<p\le x}1/p\le2/m$ dès que $\log k\ge10$.

Par conséquent, une couverture complète à ces paramètres exige, pour
le \emph{même} choix des petites classes, $|S|\le(24/m)N/\log z$.
Si, pour un seul $m\ge2$, tous les choix donnent strictement plus de
survivants pour tous les $k$ assez grands, alors
\[
      Y(t)\ll_m t^{\,2m/(m+1)}\qquad(t\to\infty).
\]
On prend $k=\lceil t^{1/(m+1)}\rceil$ et on utilise la monotonie de
$Y$. Il suffirait en particulier d'une constante positive arbitraire
dans une borne inférieure uniforme $|S|\ge c z^2/\log z$ aux carrés,
puis de prendre $m>24/c$.

On peut éviter de poser cette hypothèse plus forte. Appelons ici un
profil complétable un choix des classes jusqu'à $k^{m+1}$ qui recouvre
$[1,k^{2m}]$ et qui, au stade $k^m$, laisse au plus
$(24/m)k^{2m}/\log(k^m)$ survivants. Alors les deux assertions suivantes
sont équivalentes.
\begin{enumerate}
\item Il existe $\alpha<2$ tel que $Y(t)\ll t^\alpha$ pour tout $t$
assez grand.
\item Il existe un entier $m\ge2$ tel que, pour tout $k$ assez grand,
aucun de ces profils ne soit complétable.
\end{enumerate}
Pour le premier sens, choisir $m$ avec $(m+1)\alpha<2m$. La borne de
croissance interdit alors de couvrir $k^{2m}$ positions au seuil
$k^{m+1}$. Réciproquement, la majoration des pertes force toute
couverture à satisfaire la condition de faible effectif. L'interpolation
entre les puissances donne donc l'exposant $2m/(m+1)$. Ces implications
ne décident pas laquelle des configurations peut effectivement exister.

Cette hypothèse uniforme rencontre un obstacle précis. Sous l'hypothèse
d'une infinité de zéros de Siegel, le corollaire 1 de Granville
\cite{granville}, appliqué à $v=2$ où $f(2)=0$, fournit des intervalles
de longueur $z^2$ contenant $o(z^2/\log z)$ survivants. La translation
de l'intervalle définit les petites classes. Pour atteindre un seuil
$k^m$ avec $m$ fixé, on arrondit $z$ vers le haut. L'ajout des
premiers retire des positions, et le rallongement de l'intervalle
ajoute au plus $O_m(z^{2-1/m})=o(z^2/\log z)$ positions.
Ces exemples conditionnels concernent le nombre initial de survivants.
Il reste à savoir si les classes plus grandes peuvent effectivement
les couvrir tous.

\paragraph{Représentant commun des petites classes.}
La condition d'effectif ci-dessus exclut inconditionnellement une
famille à l'échelle des puissances. Supposons $z\ge512$ et que les
restes jusqu'à $z$ coïncident avec un même entier $b\in[1,z^2]$.
Du côté le plus long de l'intervalle à partir de $b$, chaque distance
première $z<q\le\lfloor z^2/2\rfloor$ fournit un survivant distinct.
En effet $n=b+q$ ou $n=b-q$, et aucun premier $p\le z$ ne divise
$n-b$. En notant $h=\lfloor z^2/2\rfloor$, les estimations
élémentaires de Chebyshev donnent
\[
 |S|\ge\pi(h)-\pi(z)\ge\frac{z^2}{64\log z}.
\]
On utilise $\pi(h)\ge h/(4\log h)$, $4h\ge z^2$,
$\log h\le2\log z$ et $\pi(z)\le8z/\log z$.
Pour $z=k^m$ et $m>1536$, ce minorant dépasse strictement
$(24/m)z^2/\log z$. Aucun choix des classes restantes jusqu'à
$k^{m+1}$ ne peut couvrir l'intervalle pour $k$ suffisamment grand.
Le représentant commun intervient dans l'hypothèse. Le théorème chinois
seul ne donne pas un représentant qui appartienne à l'intervalle.

On peut aussi changer certaines classes. Pour deux choix $a,c$, posons
$E=\{p\le z:\ p\ {\rm premier},\ a_p\not\equiv c_p\bmod p\}$
et $M=\sum_{p\in E}1/p$. Si $z\ge1024$, alors
\[
 \bigl|\,|S_z(a)|-|S_z(c)|\,\bigr|
 \le 68\,\frac{z^2}{\log z}\,M ,
\]
où $S_z(a)$ désigne les survivants dans $[1,z^2]$.
En effet, chaque entier de $S_z(a)\setminus S_z(c)$ appartient à la
nouvelle classe d'au moins un $p\in E$. Dans une telle classe on écrit
$n=pj+r$, avec $0\le r<p$. Les indices $0\le j<J=\lfloor z^2/p\rfloor+1$
évitent une classe pour chaque premier $q\le z$ distinct de $p$, par
inversion de $p$ modulo $q$. Chaque indice évite aussi au moins l'une
des classes $0,1$ modulo $p$. Les deux cribles ainsi complétés ont
chacun au plus $17J/\log J$ survivants, car $J\ge z\ge1024$ et
$\lfloor J^{1/8}\rfloor\le z$. Comme $J\le2z^2/p$, la somme sur $E$
donne l'inégalité dans un sens, puis on échange $a,c$.

Pour $m=3072$ fixé, $z=k^m$ et $k$ suffisamment grand, un recouvrement
complet jusqu'à $k^{m+1}$ doit donc satisfaire, pour chaque $1\le b\le z^2$,
\[
 \sum_{\substack{p\le z\ {\rm premier}\\a_p\not\equiv b\bmod p}}\frac1p
 \ge \frac{1/64-24/3072}{68}=\frac1{8704}.
\]
Cette condition permet un nombre variable de classes exceptionnelles.

On peut isoler les écarts sur les petits premiers. Fixons un entier
$D\ge1$ tel que $a_p\not\equiv b\bmod p$ pour chaque premier $p\mid D$,
$p\le z$. Prenons comme classes auxiliaires $a_p$ pour $p\mid D$ et
$b\bmod p$ ailleurs. Du côté le plus long de $b$, tous les entiers
$b\pm Dq$, pour $z<q\le h=\lfloor z^2/(2D)\rfloor$ premier, survivent.
En effet, si $p\mid D$ ils sont congrus à $b$ et, sinon, $p\nmid Dq$.
Si $z\ge512D$, alors $z\le h$, $z^2\le4Dh$ et $\log h\le2\log z$.
Les mêmes bornes de Chebyshev donnent
$D(\pi(h)-\pi(z))\ge z^2/(64\log z)$. Appliquons la stabilité aux
classes auxiliaires et aux classes du recouvrement. Pour $m,D$ fixés
et $k$ suffisamment grand, il faut donc
\[
 \frac1{64D}\le\frac{24}{m}
 +68\sum_{\substack{p\le z\ {\rm premier},\ p\nmid D\\
                       a_p\not\equiv b\bmod p}}\frac1p .
\]
Par exemple, si la parité diffère de celle de $b$, prendre $D=2$ et
$m=6144$ impose une somme d'au moins $1/17408$ sur les premiers impairs.

Soient $a>0$, $K>0$, $C>0$ fixés et
$H=\lfloor x(\log x)^a\rfloor$. Supposons qu'il existe pour une
infinité de $x$ des classes recouvrant $[1,H]$ et un entier $b$ avec
$|b|\le CH$ tel que $a_p\equiv b\bmod p$ pour tout premier
$p\le x/(\log x)^K$. Alors
\[
                 a\le \frac{K}{\sqrt e-1}.
\]
En effet, comptons les $n\in[1,H]$ tels que $|n-b|$ soit un premier
supérieur à $x$. Leur nombre est $(1+o(1))H/\log x$, uniformément pour
ces $b$. Cela suit du théorème des nombres premiers aux extrémités
des deux intervalles de distances positives, toutes comprises entre
$0$ et $(C+1)H$. Plus précisément $\pi(t)=t/\log H+o(H/\log H)$
uniformément sur cette plage, en séparant $t\le H/(\log H)^2$.
Les distances au plus $x$ coûtent au plus $2\pi(x)=o(H/\log H)$.

Aucun petit premier dans l'hypothèse ne couvre un tel $n$. Pour
$x/(\log x)^K<p\le x$, la progression $n\equiv a_p\bmod p$
a $M=H/p+O(1)$ termes. Sur ses indices, les conditions
$r\nmid n-b$ pour $r\le\sqrt M$ interdisent chacune une classe,
puisque $r<p$. Elles sont nécessaires pour tout entier compté.
Le crible supérieur au paramètre $2$ donne au plus
$(2+o(1))H/[p\log(H/p)]$ tels entiers. L'erreur est uniforme puisque
$M\ge (1+o(1))(\log x)^a\to\infty$. Une sommation partielle du théorème
des nombres premiers donne
\[
 \sum_{x/L^K<p\le x}\frac1{p\log(H/p)}
 \sim \int_{x/L^K}^x\frac{dt}{t\log t\log(H/t)}
 \sim \frac{\log(1+K/a)}{L}.
\]
La dernière équivalence utilise $\log t\sim L$ et la substitution
$v=\log(H/t)$. La couverture exige donc
$1\le 2\log(1+K/a)$, ce qui prouve l'affirmation.
En particulier, une construction de couverture à l'échelle
$x(\log x)^{1+o(1)}$ ne peut fixer tous ces petits restes à zéro
avec $K<\sqrt e-1$. On applique le résultat sur un sous-intervalle
$[1,\lfloor x(\log x)^a\rfloor]$ avec un $a<1$ fixé suffisamment
proche de $1$, en gardant $b=0$.

\paragraph{Portée.}
La substitution $k=\pi(x)$ dans le théorème 1.1 de \cite{survivors}
donne directement $Y(x)\ll x^2/[(\log x)^2(\log\log x)^2]$.
La proposition utilise aussi le comptage de survivants du théorème 1.2
et la répartition des premiers entre $z$ et $x$. Elle améliore
asymptotiquement cette conséquence directe. L'idée de sommer les
pertes classe par classe figure dans \cite[section 2.3]{bft}.
La note n'établit
ni une priorité bibliographique ni un minorant comparable.
Les constructions de Ford et al. \cite{fgkmt} et du crible incliné
\cite{tilted} demeurent d'une échelle différente. La seconde établit
$Y(x)\gg x\log x/\log_3 x$, où $\log_j$ désigne le logarithme itéré
$j$ fois. Le noyau Lean 4.34.1 a vérifié la majoration avec un
$\eta>0$ existentiel dans \texttt{erdos687\_log\_saving}, en utilisant le
théorème formalisé de la source. Son audit donne seulement
\texttt{propext}, \texttt{Classical.choice} et \texttt{Quot.sound}.
Le fichier \texttt{BorneExplicite687.lean} certifie en outre
$\eta=1/100\,000$ avec les mêmes trois axiomes. Le coefficient
$\exp(\gamma)/100$ de la source donne un coefficient de réserve
$\exp(\gamma)/2400$. Pour $d=1/200\,000$, la perte $3\cdot35d$
est strictement inférieure à cette réserve, et la comparaison des
longueurs donne $\eta=2d$. La chaîne native plus précise vérifie
également $\eta=1/200$. Le fichier \texttt{StockFinNatif687.lean}
conserve la marge $e^\gamma/40$ de la référence, puis un coefficient
de réserve $3267e^\gamma/800000$. Le fichier \texttt{CribleFinNatif687.lean}
borne chaque progression avec le coefficient $9/4$. À $d=1/400$,
le majorant harmonique de la bande a pour coefficient $(21/20)d$.
L'inégalité stricte
\[
  \frac{21}{20}\frac1{400}\frac94 < \frac{3267e^\gamma}{800000}
\]
laisse un survivant pour chaque choix de classes.
Le noyau Lean 4.34.1 a vérifié \texttt{TransfertFin687.gain\_one\_two\_hundred}
avec les trois axiomes indiqués ci-dessus.
Si $f(H)$ désigne le plus petit seuil de premiers permettant de couvrir
$[1,H]$, la même borne donne
\[
       f(H)\gg \sqrt H\,(\log H)^{401/400}.
\]
En effet $H\le Y(f(H))$ et la majoration élémentaire utilisée dans cette
comparaison donne $H\le f(H)^3$ pour $H$ assez grand. Donc
$\log H\le3\log f(H)$. Il suffit de prendre la racine carrée dans la
borne supérieure de $Y(f(H))$. Le noyau natif vérifie cette implication
dans \texttt{InverseFin687.cutoff\_lower}.

La preuve conserve aussi une marge si l'on ajoute des classes.
Pour $d=1/400$ et pour la même longueur $N=\mathrm{length}(d,x)$
utilisée dans le transfert, posons $L=\log x$ et
$W=N\log L/L^2$. Il existe deux constantes $\kappa,c'>0$ telles que,
pour tout $x$ suffisamment grand, après le criblage par tous les
premiers jusqu'à $x$ puis par un ensemble arbitraire de premiers $E$
au-dessus de $x$, au moins $c'W$ entiers survivent dès que
$|E|\le\kappa x\log L/L$. En effet, si $m$ est la différence strictement
positive entre le coefficient de réserve et $(21/20)d(9/4)$,
la première étape conserve au moins $mW$. Une classe supplémentaire coûte
au plus $35N/(xL)$. Les choix $\kappa=m/70$ et $c'=m/2$ conviennent.
Le noyau Lean 4.34.1 vérifie \texttt{SurvieFine687.additional\_prime\_budget}.

\paragraph{Agrandissement du seuil.}
L'application de cette réserve à tous les premiers d'une bande donne
directement une majoration plus forte de la fonction demandée :
\[
 Y(t)\ll \frac{t^2}{(\log t)^{401/200}(\log\log t)^2}.
\]
En effet, fixons $A=\max(1,16/\kappa)$ et prenons
$s=\lceil At/\log\log t\rceil$. Finalement $s\le t$,
$\log t/2\le\log s\le\log t$ et
$\log\log t/2\le\log\log s$. La majoration élémentaire
$\pi(t)\le8t/\log t$ donne
\[
   \#\{p:s<p\le t\}\le \pi(t)
       \le \kappa s\,\frac{\log\log s}{\log s}.
\]
La réserve démontrée au seuil $s$ permet ainsi d'ajouter \emph{tous}
les premiers manquants jusqu'à $t$. Un entier de l'intervalle de
longueur $N=\mathrm{length}(1/400,s)$ survit à tout choix des classes.
Donc $Y(t)\le N\ll s^2/(\log s)^{401/200}$. Comme
$s\le2At/\log\log t$ pour $t$ assez grand, la borne annoncée suit.
Le noyau Lean 4.34.1 a vérifié \texttt{BorneIteree687.upper} avec
\texttt{propext}, \texttt{Classical.choice} et \texttt{Quot.sound}.
Le facteur supplémentaire tend vers l'infini. La question d'un
minorant comparable reste ouverte.

\begin{thebibliography}{9}
\bibitem{survivors} OpenAI, \emph{A quadratic bound for Jacobsthal's
function}, prépublication, 25 septembre 2026, théorèmes 1.1 et 1.2.
\href{https://github.com/openai/math/blob/adc7f1241b42e322a6451854ab7e4b4c146bf78a/preprints/A-quadratic-bound-for-Jacobsthals-function-September-25-2026/paper.pdf}{Manuscrit et preuve.}
\bibitem{granville} A. Granville, \emph{Sieving intervals and Siegel zeros},
Acta Arithmetica 205 (2022), 1--19, formule (1), crible
linéaire de Jurkat et Richert.
\href{https://arxiv.org/abs/2010.01211}{Texte.}
\bibitem{bft} W. Banks, K. Ford et T. Tao, \emph{Large prime gaps and
probabilistic models}, section 2.3.
\href{https://ford126.web.illinois.edu/wwwpapers/gaps-model.pdf}{Texte.}
\bibitem{fgkmt} K. Ford, B. Green, S. Konyagin, J. Maynard et T. Tao,
\emph{Long gaps between primes}, Journal of the American Mathematical
Society 31 (2018), 65--105.
\href{https://arxiv.org/abs/1412.5029}{Texte.}
\bibitem{tilted} GPT-5.6 Sol, manuscrit diffusé par DottedCalculator,
\emph{A tilted residue-class construction for long prime-free intervals},
25 août 2026, théorème (1.2).
\href{https://github.com/DottedCalculator/ai-math/blob/main/Erdos_4_GPT_5.6_Sol.pdf}{Manuscrit.}
\href{https://github.com/plby/lean-proofs/blob/main/src/latest/ErdosProblems/Erdos4Tilted.lean}{Formalisation de B. Alexeev.}
\end{thebibliography}
\end{document}
